\subsection{Normed spaces on a valued division ring}

Recall (GT, IX, § 3.2) that an absolute value on a division ring K is a mapping \(\xi \mapsto |\xi|\) of K in \(\mathbf{R}_{+}\) , such that \(|\xi| = 0\) if, and only if, \(\xi = 0\) , and that \(|\xi \eta| = |\xi| \cdot |\eta|\) , and \(|\xi + \eta| \leqslant |\xi| + |\eta|\) ; an absolute value defines a distance \(|\xi - \eta|\) on K, and hence a Hausdorff topology compatible with the division ring structure of K. If \(|\xi| = 1\) for all \(\xi \neq 0\) , the absolute value is called improper, and the topology that it defines on K is the discrete topology; if, on the other hand, there exists \(\alpha \neq 0\) in K such that \(|\alpha| \neq 1\) , then there exists \(\beta \neq 0\) in K such that \(|\beta| < 1\) (it is sufficient to take \(\beta = \alpha\) or \(\beta = \alpha^{-1}\) ), and the sequence \((\beta^n)_{n \geqslant 1}\) converges to 0, thus the topology of K is not discrete.

We recall on the other hand (GT, IX, § 3.3) that if E is a vector space on a non-discrete valued division ring K then a norm on E is a mapping \(x \mapsto \|x\|\) of E in \(\mathbf{R}_{+}\) , such that \(\|x\| = 0\) if, and only if, \(x = 0\) , and such that \(\| \lambda x \| = |\lambda| \cdot \| x\|\) for every scalar \(\lambda \in \mathbf{K}\) , and \(\|x + y\| \leqslant \|x\| + \|y\|\) . A distance \(\|x - y\|\) , is defined on E by the norm, and hence a topology that is compatible with the vector space structure of E (loc. cit.). Unless the contrary is expressly stated, a normed space is considered in terms of the structure of the topological vector space defined by its norm. The normed spaces are among the most important of topological vector spaces.

It is known (GT, IX, § 3.3) that two distinct norms on E can define the same topology on E; for this it is necessary and sufficient that the two norms be equivalent (loc. cit.). The structure of normed spaces is thus richer than the structure of topological vector spaces; if E and F are two normed spaces one must be careful to distinguish between the idea of isomorphism of the normed space structure of E with that of F, and the idea of isomorphism of the topological vector space structure of E with that of F.

Example. --- Let I be an arbitrary set of indices; it is known (GT, X, § 3.2) that a norm \(\|x\|\) can be defined, on the set of bounded mappings \(x = (\xi_{\mathfrak{t}})\) of I in K, \(\mathcal{B}(\mathrm{I};\mathrm{K})\) (also written \(\mathcal{B}_{\mathrm{K}}(\mathrm{I})\) or \(\ell_{\mathrm{K}}^{\infty}(\mathrm{I})\) ), by \(\|x\| = \sup_{\mathfrak{t}\in \mathrm{I}}|\xi_{\mathfrak{t}}|\) . When I is a topological space, the set of bounded, continuous mappings of I in K is a closed subspace of the space \(\mathcal{B}(\mathrm{I};\mathrm{K})\) (GT, X, § 3.1, cor. 2). Another subspace of \(\mathcal{B}(\mathrm{I};\mathrm{K})\) is the set \(\ell_{\mathrm{K}}^{1}(\mathrm{I})\) of absolutely summable families \(x = (\xi_{\mathfrak{t}})\) (GT, X, § 3.6); we can define on this subspace another norm \(\|x\|_1 = \sum_{\mathfrak{t}\in \mathrm{I}}|\xi_{\mathfrak{t}}|\) , that in general is not equivalent to the norm \(\|x\| = \sup_{\mathfrak{t}\in \mathrm{I}}|\xi_{\mathfrak{t}}|\) (I, p.~23, exerc. 6); when considering \(\ell_{\mathrm{K}}^{1}(\mathrm{I})\) as a normed space, without specifying its norm, it is always the norm \(\|x\|\) , that is meant. We write \(\mathcal{B}(\mathrm{I})\) and \(\ell^1 (\mathrm{I})\) in place of \(\mathcal{B}(\mathrm{I};\mathbf{R})\) and \(\ell_{\mathbf{R}}^{1}(\mathrm{I})\) .

